%! AP Statistics — Unit 1, Lesson 1.1 — Homework
\documentclass[10pt]{article}
\usepackage{apstats-article}
\usepackage{apstats-boxes}

%=============================================================================
\begin{document}

\pageheader{Unit 1, Lesson 1.1}{Homework}
\namedateperiod

\vspace{0.1in}

\begin{tcolorbox}[colback=sky, colframe=navy, boxrule=0.8pt, arc=2mm,
    left=3mm, right=3mm, top=1.5mm, bottom=1.5mm]
\small For each news blurb, identify the \textbf{individual}, \textbf{population},
\textbf{sample}, and \textbf{variable(s)}. Write in complete phrases, not single words.
If no sample size is stated, write \textit{not given}.
\end{tcolorbox}

\vspace{0.14in}

% ════════════════════════════════════════════════════════════════════════════
% BLURB 1
% ════════════════════════════════════════════════════════════════════════════
\begin{blurbbox}[Study 1]{navylight}
\begin{headlinebox}{sky}
\small\textit{``Researchers surveyed 1{,}800 teenagers across 12 countries and
found that 67\% felt social media had a negative effect on their self-image.''}
\end{headlinebox}

\vspace{8pt}
\componenttable
\end{blurbbox}

\vspace{0.14in}

% ════════════════════════════════════════════════════════════════════════════
% BLURB 2
% ════════════════════════════════════════════════════════════════════════════
\begin{blurbbox}[Study 2]{greenacc}
\begin{headlinebox}{greenbg}
\small\textit{``A study of 250 WNBA players found that athletes who practiced
yoga at least twice a week reported significantly fewer injuries during the season
compared to those who did not.''}
\end{headlinebox}

\vspace{8pt}
\componenttable
\end{blurbbox}

\vspace{0.14in}

% ════════════════════════════════════════════════════════════════════════════
% BLURB 3
% ════════════════════════════════════════════════════════════════════════════
\begin{blurbbox}[Study 3]{redacc}
\begin{headlinebox}{redbg}
\small\textit{``Scientists collected soil samples from 45 farms in Iowa and
found that fields using cover crops contained 30\% more earthworms per square
meter than fields without cover crops.''}
\end{headlinebox}

\vspace{8pt}
\componenttable
\end{blurbbox}

\vspace{0.14in}

% ── REFLECTION ───────────────────────────────────────────────────────────────
\begin{reflectionbox}
\small Choose \textbf{one} of the three studies above and answer:
\textit{Why might the results of this study not apply to everyone?}
What would make you more --- or less --- confident in the findings?

\vspace{6pt}
\textcolor{slate}{\scriptsize Study \#:\;\blank{0.6cm}}

\vspace{4pt}
\writeline\\[1pt]\writeline\\[1pt]\writeline\\[1pt]\writeline
\end{reflectionbox}

\vspace{0.14in}

% ── EXTENSION ────────────────────────────────────────────────────────────────
\begin{extensionbox}
\small

\textbf{Part A --- U.S.\ Census}\quad
The U.S.\ government conducts a census of every person in the country every
10 years. Research (or reason through) the following and write a few sentences
for each:

\vspace{5pt}
\begin{enumerate}[leftmargin=1.5em, itemsep=8pt]
  \item What does it cost --- in money and time --- to conduct the U.S.\ Census?\\[4pt]
        \writeline\\[1pt]\writeline

  \item Name one group of people who are difficult to count accurately and explain why.\\[4pt]
        \writeline\\[1pt]\writeline

  \item Why do we bother with a census at all instead of just using a large sample?\\[4pt]
        \writeline\\[1pt]\writeline
\end{enumerate}

\vspace{8pt}
\textbf{Part B --- Preview: Categorical vs.\ Quantitative}\\[4pt]
Look back at the variables you identified in Studies 1--3 above.
For each variable, decide: does it produce a \emph{number} you can do arithmetic
with, or does it put individuals into \emph{categories}?

\vspace{6pt}
\renewcommand{\arraystretch}{1.8}
\begin{tabularx}{\linewidth}{@{} >{\bfseries}p{0.08\linewidth}
                                    >{\raggedright\arraybackslash}p{0.44\linewidth}
                                    X @{}}
& \textbf{Variable} & \textbf{Number or category?} \\
\hline
S1 & \blank{7cm} & \blank{4cm} \\
S2 & \blank{7cm} & \blank{4cm} \\
S3 & \blank{7cm} & \blank{4cm} \\
\end{tabularx}

\vspace{4pt}
\textcolor{slate}{\scriptsize We will formalize this distinction in Lesson 1.2.}
\end{extensionbox}

\end{document}
